\begin{proposition}[Irreducible fibres of torus fibrations]\hypertarget{tag-02WJ}{}\label{night:P4.12}\label{stmt:P4.12}
\StatementRecord{night:P4.12}{proved}{02WJ}{\SeriesRef{night:P4.6}}

Let $Y$ be a connected smooth compact complex threefold with
$b_1(Y)=b_2(Y)=0$, and let $f:Y\to B$ be a proper surjective
holomorphic map with connected fibres to a smooth compact curve.
Assume that the smooth general fibre is a complex two-torus, and let
$f^{-1}(b)=mS$ have irreducible reduced support. Then
\[
 \omega_S^{\otimes m}\simeq\OO_S,\qquad
 \chi(\OO_S)=0,\qquad h^1(\OO_S)\in\{1,2\}.
\]
Every normal such support is smooth. For a singular, necessarily
nonnormal support, a minimal resolved model of its normalization
has Kodaira dimension $-\infty$.
\end{proposition}
\begin{proof}
\noindent\textbf{Step 1. Canonical torsion.}
Choose a disc $\Delta$ with no other critical value, with coordinate
$t$ vanishing at $b$, and write $U=f^{-1}(\Delta)$. The coherent
direct image $f_*\omega_U$ is torsion-free of rank one: nonzero
base functions act injectively on sections, and a nearby torus has
one-dimensional canonical space. It is therefore a line bundle
on the disc. A generator gives a section $v$ of $\omega_U$,
nonzero and nowhere vanishing on every noncentral fibre. Its
divisor is $aS$ for an integer $a\ge0$. If $a\ge m$, then $v/t$
would be holomorphic, contradicting primitivity of the generator
at zero. We have $0\le a<m$ and
\[
 \omega_U\simeq\OO_U(aS),\qquad\OO_U(mS)\simeq\OO_U.
\]
Adjunction gives $\omega_S\simeq\OO_U((a+1)S)|_S$, proving
$\omega_S^{\otimes m}\simeq\OO_S$. The ambient
Riemann--Roch computation in \ref{night:P4.6}, applied to $S$, gives
$\chi(\OO_S)=0$. Reduced irreducibility gives $h^0(\OO_S)=1$.
A nonzero section of a torsion line bundle has a power equal to
a nonzero constant, and hence is nowhere zero; its space of
sections consequently has dimension at most one. Serre duality
on the Gorenstein surface $S$ gives
$h^2(\OO_S)=h^0(\omega_S)\in\{0,1\}$, proving the assertion
about $h^1$.

\noindent\textbf{Step 2. Exceptional curves and discrepancies.}
For a minimal resolution $\pi:T\to S$ of a normal Gorenstein
surface write $K_T=\pi^*K_S+\sum a_jE_j$. Its exceptional
intersection matrix $M$ has nonnegative off-diagonal entries
and is negative definite. Indeed,
pulling back a generic holomorphic function vanishing at a
contracted point gives $H+\sum m_jE_j$, with $m_j>0$ and $H$
the strict transform. Intersecting gives $Mm\le0$, strictly in
at least one entry of each connected block. After scaling by
$\operatorname{diag}(m_j)$, the negative quadratic form is a
sum of nonnegative terms $c_{ij}(z_i-z_j)^2$ and $d_i z_i^2$,
with connected support and some $d_i>0$ in every block. It is
positive definite. This is the quadratic identity of
\cite[\S I, unnumbered assertion and proof, p.~6]{Mumford1961}.

Adjunction gives
$K_T\cdot E_j=2p_a(E_j)-2-E_j^2\ge0$: the sole negative
possibility is a smooth rational $(-1)$-curve, excluded by
minimality of the resolution. The inverse of a negative definite
matrix with nonnegative off-diagonal entries has nonpositive
entries. Indeed, put $Q=-M$ and choose $c$ larger than all its
eigenvalues and diagonal entries; then
$Q^{-1}=c^{-1}\sum_{k\ge0}(I-Q/c)^k$ is entrywise nonnegative.
The discrepancy equations $Ma=(K_T\cdot E_j)_j$ imply
$a_j\le0$. Hence $K_T=\pi^*K_S-E$ with $E$ an effective
exceptional integral divisor.

Put $p=\dim H^0(S,R^1\pi_*\OO_T)$. Leray, with this sheaf
supported at finitely many points, gives
\begin{equation}\hypertarget{tag-01OD}{}\label{eq:P4chi-resolution}
 \chi(\OO_T)=\chi(\OO_S)-p=-p.
\end{equation}
For an effective exceptional cycle $Z$, put $\mathcal F=I_Z$.
Fix $s\in S$ and its maximal ideal $\mathfrak m_s$. Each coherent
quotient $\mathcal F/\mathfrak m_s^k\mathcal F$ is supported on a
compact infinitesimal fibre of dimension at most one. Its cohomology
vanishes in degrees at least two. To see this, filter by the nilradical
and reduce to coherent sheaves on the reduced compact curve. Removing
point-supported torsion and pulling back to its finite normalization
leaves vector bundles on compact Riemann surfaces, up to kernels and
cokernels supported at finitely many points. Their Dolbeault resolutions
have length one, while point-supported sheaves have no higher
cohomology. The finite filtration therefore gives the asserted vanishing
on every infinitesimal fibre.
Analytic formal functions and coherence
\cite[\S6, Haupts\"atze~I and~IIa, p.~59]{Grauert1960D} identify
the completed stalk of $R^2\pi_*\mathcal F$ with the inverse limit of
these zero groups. Nakayama's lemma gives $R^2\pi_*I_Z=0$.
The exact sequence of $Z$ gives a surjection
$R^1\pi_*\OO_T\to R^1\pi_*\OO_Z$, and therefore
\begin{equation}\hypertarget{tag-01OE}{}\label{eq:P4exceptional-genus}
 p\ge h^1(Z,\OO_Z).
\end{equation}

\noindent\textbf{Step 3. Exclusion of normal singular supports.}
Suppose first that $E=0$. Then $K_T$ is torsion and the
classification \cite[Chapter~VI, Theorem~1.1, p.~243]{BHPV2004E} gives Kodaira dimension zero and nonnegative
holomorphic Euler characteristic. Equation
\eqref{eq:P4chi-resolution} forces $p=0$ and $\chi(\OO_T)=0$.
The remaining surfaces are tori, bielliptic surfaces and Kodaira
surfaces. They have no rational curves: a map from $\Pp^1$ to
a torus is constant, bielliptic surfaces have finite unramified
torus covers, and a primary Kodaira surface has an elliptic
fibration with smooth elliptic reduced fibres over an elliptic
curve. A map from $\Pp^1$ to a secondary Kodaira surface lifts
to its finite unramified primary Kodaira cover. On a torsion
canonical surface a negative curve is a smooth rational
$(-2)$-curve by adjunction. Thus there are no exceptional
curves, and $S$ is smooth.

Suppose next that $E\ne0$. Since $K_T+E$ is torsion, all
positive plurigenera vanish, so $T$ is rational, ruled, or a
blowup of a class-VII surface. Rationality contradicts
\eqref{eq:P4chi-resolution}, since $\chi(\OO_T)=1$. In class
VII that equation gives $p=0$, whereas
\[
 \chi(\OO_E)=-E\cdot(E+K_T)/2=0.
\]
The nonempty effective compact divisor $E$ has $h^0(\OO_E)\ge1$,
so $h^1(\OO_E)\ge1$, contradicting
\eqref{eq:P4exceptional-genus}.

Finally suppose $T$ is ruled, after point blowups, over a curve
$C$ of genus $g\ge1$. Equation \eqref{eq:P4chi-resolution}
gives $p=g-1$. A general ruling fibre $F$ has $K_T\cdot F=-2$,
so $E\cdot F=2$. An exceptional component $E_j$ therefore
maps onto $C$. Its normalization has genus at least $g$ by
Riemann--Hurwitz, so $h^1(\OO_{E_j})\ge g$. Taking $Z=E_j$
in \eqref{eq:P4exceptional-genus} gives $p\ge g$, a contradiction.

\noindent\textbf{Step 4. The resolved normalization of a nonnormal support.}
Write $\nu:N\to S$ for the finite normalization and $\pi:T\to N$
for its minimal resolution. The residue comparison of dualizing
sheaves gives a pole divisor along the conductor; its strict
transform $D$ on $T$ is nonzero and effective. Pull back a local generator of $\omega_S$ as a meromorphic
two-form. Its canonical divisor has the form $-D+\sum a_jE_j$.
The equations for the exceptional coefficients have right side
$(K_T+D)\cdot E_j\ge0$, because both $K_T\cdot E_j$ and
$D\cdot E_j$ are nonnegative. The same inverse-matrix argument
gives $a_j\le0$. Taking the $m$th power to trivialize the pulled
back $\omega_S$ consequently gives
\[
 mK_T+D'\sim0
\]
for a nonzero effective integral divisor $D'$. This is also the
conductor formula of \cite[Lemma~4.6, p.~7 of the attached-source edition]{CDP20}. Any
positive pluricanonical section on $T$, multiplied by the
appropriate effective section and raised to a power, would give
a nonzero holomorphic function on compact $T$ with a zero.
It follows that $\kappa(T)=-\infty$. Contraction of exceptional
curves of the first kind preserves the plurigenera, so the minimal
model has the same Kodaira dimension.

\end{proof}